\section{Synthetic and mechanism studies}

Synthetic and controlled mechanism studies are organized by the question they
test rather than by dataset. Because the latent biological and acquisition
factors are known in synthetic generators, recovery can be measured rather than
inferred indirectly. All synthetic claims are kept strictly separate from
pathology claims.

\subsection{Identifiability}

Synthetic crossed-factor identifiability studies ask when ordinary observations
and paired interventions jointly control recovery of latent biological and
acquisition factors, covering observational sample size, unique paired-anchor
count, factor overlap, linear versus nonlinear mixing, and seed. Crossed-target
scanner-prototype factorization and unseen-identity representation-geometry
studies examine factor recovery under crossed targets and held-out identities.
Finite-sample whitening identifiability audits examine when whitening aids or
confounds recovery.

\subsection{Factor allocation and capacity}

Paired-acquisition dimensionality $\times$ cross-covariance factorials (450
cells) and capacity-matched campaigns (175 fits) probe the separation frontier
of biological and acquisition branches. The biological-bottleneck capacity
allocation factorial (status
\texttt{\seqsplit{complete\_capacity\_gain\_with\_scanner\_tradeoff}}) showed a capacity
gain with a scanner tradeoff. Routed paired-consensus bottleneck and
paired-consensus linear-anchor studies test routed and consensus supervision.
The real paired-scanner bottleneck-allocation validation
(\texttt{\seqsplit{complete\_mixed\_real\_paired\_scanner\_allocation\_effects}}) transferred
the capacity question to real paired-scanner data: B64 improved retrieval and
increased scanner recoverability without a category gain.

\subsection{Task sufficiency}

Task-defined biological-sufficiency benchmarks and task-benchmark
instrument-power audits measure how much capacity is actually used by a task and
whether the evaluation instrument has power to detect differences. Minimal and
whitened biological-bottleneck studies and residual-probe nonlinear-capacity
calibration examine how representation geometry responds to constrained
capacity.

\subsection{Synthetic-to-real transport}

Synthetic accessibility effects do not automatically transport to real pathology
features. The corrected fixed-estimand adjudication established that the
synthetic width effect did not produce a corrected real category gain:
\texttt{\seqsplit{synthetic\_accessibility\_effect\_transported}} is false because B64 did
not improve corrected five-category canine balanced accuracy. Retrieval gain is
never counted as transport of biological accessibility.

\subsection{Interpretation boundary}

Synthetic results demonstrate mechanism, identifiability, and allocation
behavior under controlled generators. They are not pathology evidence, do not
establish clinical utility, and are never used to inflate a real-data claim
beyond what the corrected real-data artifacts support.
